% GENERATED FILE. Edit canonical lessons, then run npm run generate:books.
% canonical-source-hash: fnv1a64:f6e52140d8cf95a1
% canonical-lessons: KA-C125-idli, KA-C125-dose, KA-C125-vade, KA-C125-uppittu, KA-C125-ragi

\chapter{Idli, Dosa, Vada, Upma, Finger millet}
\label{ch:ka-idli-dosa-vada-upma-finger-millet}
\hlchaptermodality{125}

\begin{chapteropening}
\textbf{By the end of this chapter:} \emph{I can say an idli, a dosa, a vada, upma and finger millet.}

Complete the last lesson of chapter 125: I can say an idli, a dosa, a vada, upma and finger millet.
\end{chapteropening}

\section[\texorpdfstring{\kn{ಇಡ್ಲಿ} (iḍli)}{ಇಡ್ಲಿ (iḍli)}]{\kn{ಇಡ್ಲಿ} (iḍli) — an idli}

\label{lesson:KA-C125-idli}



\begin{quote}
Before the new one: say the Kannada for garlic, then the Kannada for honey.
\end{quote}

\subsection*{You'll want to know: \kn{ಇಡ್ಲಿ}}
\textbf{\kn{ಇಡ್ಲಿ}} — \emph{iḍli} — \textquotedblleft{}an idli\textquotedblright{}.

\begin{grammarlens}[title={What you've built}]
One more word for this chapter.
\end{grammarlens}

\subsection*{Your turn}
\begin{itemize}
\raggedright
  \item \emph{Say it:} \emph{iḍli}
  \item \emph{Say it:} \emph{iḍli}, once more
  \item \emph{Say it:} say \emph{iḍli}
  \item \emph{Recall:} say the Kannada for garlic, then the Kannada for honey, then say \emph{iḍli} again
\end{itemize}

\begin{tcolorbox}[breakable,colback=teal!4,colframe=teal!35!black,title={Before you move on}]
What does \kn{ಇಡ್ಲಿ} mean? (\textquotedblleft{}An idli\textquotedblright{}.) Say it once more.
\end{tcolorbox}
\section[\texorpdfstring{\kn{ದೋಸೆ} (dōse)}{ದೋಸೆ (dōse)}]{\kn{ದೋಸೆ} (dōse) — a dosa}

\label{lesson:KA-C125-dose}



\begin{quote}
Before the new one: say the Kannada for honey, then the Kannada for an idli.
\end{quote}

\subsection*{You'll want to know: \kn{ದೋಸೆ}}
\textbf{\kn{ದೋಸೆ}} — \emph{dōse} — \textquotedblleft{}a dosa\textquotedblright{}.

\begin{grammarlens}[title={What you've built}]
One more word for this chapter.
\end{grammarlens}

\subsection*{Your turn}
\begin{itemize}
\raggedright
  \item \emph{Say it:} \emph{dōse}
  \item \emph{Say it:} \emph{dōse}, once more
  \item \emph{Say it:} say \emph{dōse}
  \item \emph{Recall:} say the Kannada for honey, then the Kannada for an idli, then say \emph{dōse} again
\end{itemize}

\begin{tcolorbox}[breakable,colback=teal!4,colframe=teal!35!black,title={Before you move on}]
What does \kn{ದೋಸೆ} mean? (\textquotedblleft{}A dosa\textquotedblright{}.) Say it once more.
\end{tcolorbox}
\section[\texorpdfstring{\kn{ವಡೆ} (vaḍe)}{ವಡೆ (vaḍe)}]{\kn{ವಡೆ} (vaḍe) — a vada}

\label{lesson:KA-C125-vade}



\begin{quote}
Before the new one: say the Kannada for an idli, then the Kannada for a dosa.
\end{quote}

\subsection*{You'll want to know: \kn{ವಡೆ}}
\textbf{\kn{ವಡೆ}} — \emph{vaḍe} — \textquotedblleft{}a vada\textquotedblright{}.

\begin{grammarlens}[title={What you've built}]
One more word for this chapter.
\end{grammarlens}

\subsection*{Your turn}
\begin{itemize}
\raggedright
  \item \emph{Say it:} \emph{vaḍe}
  \item \emph{Say it:} \emph{vaḍe}, once more
  \item \emph{Say it:} say \emph{vaḍe}
  \item \emph{Recall:} say the Kannada for an idli, then the Kannada for a dosa, then say \emph{vaḍe} again
\end{itemize}

\begin{tcolorbox}[breakable,colback=teal!4,colframe=teal!35!black,title={Before you move on}]
What does \kn{ವಡೆ} mean? (\textquotedblleft{}A vada\textquotedblright{}.) Say it once more.
\end{tcolorbox}
\section[\texorpdfstring{\kn{ಉಪ್ಪಿಟ್ಟು} (uppiṭṭu)}{ಉಪ್ಪಿಟ್ಟು (uppiṭṭu)}]{\kn{ಉಪ್ಪಿಟ್ಟು} (uppiṭṭu) — upma}

\label{lesson:KA-C125-uppittu}



\begin{quote}
Before the new one: say the Kannada for a dosa, then the Kannada for a vada.
\end{quote}

\subsection*{You'll want to know: \kn{ಉಪ್ಪಿಟ್ಟು}}
\textbf{\kn{ಉಪ್ಪಿಟ್ಟು}} — \emph{uppiṭṭu} — \textquotedblleft{}upma\textquotedblright{}.

\begin{grammarlens}[title={What you've built}]
One more word for this chapter.
\end{grammarlens}

\subsection*{Your turn}
\begin{itemize}
\raggedright
  \item \emph{Say it:} \emph{uppiṭṭu}
  \item \emph{Say it:} \emph{uppiṭṭu}, once more
  \item \emph{Say it:} say \emph{uppiṭṭu}
  \item \emph{Recall:} say the Kannada for a dosa, then the Kannada for a vada, then say \emph{uppiṭṭu} again
\end{itemize}

\begin{tcolorbox}[breakable,colback=teal!4,colframe=teal!35!black,title={Before you move on}]
What does \kn{ಉಪ್ಪಿಟ್ಟು} mean? (\textquotedblleft{}Upma\textquotedblright{}.) Say it once more.
\end{tcolorbox}
\section[\texorpdfstring{\kn{ರಾಗಿ} (rāgi)}{ರಾಗಿ (rāgi)}]{\kn{ರಾಗಿ} (rāgi) — finger millet}

\label{lesson:KA-C125-ragi}



\begin{quote}
Before the new one: say the Kannada for a vada, then the Kannada for upma.
\end{quote}

\subsection*{You'll want to know: \kn{ರಾಗಿ}}
\textbf{\kn{ರಾಗಿ}} — \emph{rāgi} — \textquotedblleft{}finger millet\textquotedblright{}.

\begin{grammarlens}[title={What you've built}]
One more word for this chapter.
\end{grammarlens}

\subsection*{Your turn}
\begin{itemize}
\raggedright
  \item \emph{Say it:} \emph{rāgi}
  \item \emph{Say it:} \emph{rāgi}, once more
  \item \emph{Say it:} say \emph{rāgi}
  \item \emph{Recall:} say the Kannada for a vada, then the Kannada for upma, then say \emph{rāgi} again
\end{itemize}

\begin{tcolorbox}[breakable,colback=teal!4,colframe=teal!35!black,title={Before you move on}]
What does \kn{ರಾಗಿ} mean? (\textquotedblleft{}Finger millet\textquotedblright{}.) Say it once more.
\end{tcolorbox}
